\documentclass[11pt]{article}
\usepackage{amsmath, amssymb}
\usepackage{geometry}
\usepackage{booktabs}
\usepackage{tabularx}
\usepackage{float}
\floatstyle{ruled}
\newfloat{algorithm}{H}{loa}
\floatname{algorithm}{Algorithm}
\usepackage{hyperref}
\geometry{margin=1in}

\renewcommand{\arraystretch}{1.3}

% Coefficient vector for tau polynomial (distinct from aggregate rating beta_i)
\newcommand{\bsc}{\mathbf{c}}

\title{Supplemental Material: Empirical Methodology}
\date{}

\begin{document}
\maketitle

\tableofcontents
\bigskip

\noindent This document describes the construction of advertiser value and externality
scores from the XNP400 dataset, the three auction mechanisms used in the
empirical simulations, the genetic algorithm used to optimize the threshold
function, and the expected penalty curves derived from the optimized threshold.
Notation follows Tables~\ref{tab:theory_notation} and~\ref{tab:empirical_notation}
in the main text.

% ─────────────────────────────────────────────────────────────────────────────
\section{Dataset and Score Construction}
\label{sec:data}
% ─────────────────────────────────────────────────────────────────────────────

\subsection{XNP400 Dataset}

We use the XNP400 dataset, a corpus of $N = 11{,}853$ posts (after filtering)
collected from the X (Twitter) platform as of February~2, 2025.
For each post $i$ we observe:

\begin{center}
\begin{tabularx}{\textwidth}{llX}
\toprule
Symbol & Source field & Description \\
\midrule
$\eta_i$ & \texttt{impression\_count} & Total lifetime impressions \\
$\alpha_i$ & \texttt{action\_count}   & Total engagements (likes, retweets, replies,
                                        quotes, bookmarks) \\
$\mu_i$   & derived                   & Age in months: $(\text{Feb 2, 2025} -
                                        \text{created\_at}_i)/30$ \\
\bottomrule
\end{tabularx}
\end{center}

\noindent
Each post is also associated with a set of Community Notes $\mathcal{J}(i)$.
Note $j \in \mathcal{J}(i)$ carries a binary direction $\pi_{ij} \in \{-1,+1\}$
and receives a set of helpfulness ratings $\mathcal{L}(i,j)$ from the community.

\subsection{Data Collection and Filtering}
\label{sec:filtering}

Posts were retrieved from the X Community Notes public data release as of
February~2, 2025.
We restrict the sample to posts that satisfy all of the following criteria:

\begin{enumerate}
    \item \textbf{At least one Community Note.}
    We include only posts that attracted at least one note.
    Posts without any Community Note carry no externality signal and are excluded.

    \item \textbf{Positive impressions.}
    Posts with zero recorded lifetime impressions are dropped, as their
    per-month scores are undefined.

    \item \textbf{Outlier trimming.}
    We remove the top and bottom 0.1\% of posts by aggregate post rating
    $\beta_i$ and by total action count $\alpha_i$.
    This trims approximately 24 posts at each tail and prevents extreme outliers
    from distorting the score distributions.

    \item \textbf{Minimum agreement score.}
    Posts with an agreement score below $-400$ (rating units per 1{,}000
    impressions) are removed.
    These posts have extraordinarily negative community assessments that are
    likely the result of viral coordinated flagging rather than genuine
    advertiser-relevant externalities.
\end{enumerate}

After applying all filters the sample contains $N = 11{,}853$ posts spanning
posts created from the platform's early period through early 2025.

\subsection{Community Notes Scoring}
\label{sec:notes}

Each rating $l \in \mathcal{L}(i,j)$ is mapped to a helpfulness score:
\begin{equation}
    \rho_{ijl} =
    \begin{cases}
        1   & \texttt{helpfulnessLevel} = \texttt{HELPFUL} \\
        0.5 & \texttt{helpfulnessLevel} = \texttt{SOMEWHAT\_HELPFUL} \\
        0   & \texttt{helpfulnessLevel} = \texttt{NOT\_HELPFUL}
    \end{cases}
\end{equation}

The note direction encodes whether the note author classified the post as misleading:
\begin{equation}
    \pi_{ij} =
    \begin{cases}
        -1 & \texttt{classification} = \texttt{MISINFORMED\_OR\_POTENTIALLY\_MISLEADING} \\
        +1 & \texttt{classification} = \texttt{NOT\_MISLEADING}
    \end{cases}
\end{equation}

The aggregate rating of post $i$ sums the signed, weighted contributions of all
notes and their ratings:
\begin{equation}
    \beta_i = \sum_{j \in \mathcal{J}(i)} \pi_{ij}
              \Bigl(1 + \sum_{l \in \mathcal{L}(i,j)} \rho_{ijl}\Bigr),
    \label{eq:beta}
\end{equation}
where the additive~$1$ inside the parentheses treats the note author as an
implicit rater who agrees with their own note.  Posts with no notes have
$\beta_i = 0$.  A post flagged by many raters as misleading has $\beta_i \ll 0$;
one rated as helpful has $\beta_i > 0$.

\subsection{Value and Externality Scores}
\label{sec:scores}

Both scores are expressed as monthly rates to control for post age,
following the E1 and V1 methods described in the main text.

\subsubsection*{Externality score (E1)}

The per-post externality (dollars per month) is
\begin{equation}
    e_i = \zeta \cdot \frac{\beta_i}{\mu_i},
    \label{eq:e_score}
\end{equation}
where $\zeta \geq 0$ is the \emph{externality cost} parameter (dollars per unit of
aggregate rating per month).  Because most posts that attract Community Notes are
flagged as misleading, $\beta_i < 0$ for the majority of posts in the sample, so
$e_i < 0$ at any positive $\zeta$.

\subsubsection*{Advertiser value score (V1)}

The per-post advertiser value (actions per month) is
\begin{equation}
    v_i = \theta \cdot \frac{\alpha_i}{\mu_i},
    \label{eq:v_score}
\end{equation}
where $\theta \geq 0$ is the \emph{action cost} parameter (dollars per engagement
action per month).  We set $\theta = 1$ throughout, so $v_i$ is measured directly
in actions per month and serves as the advertiser's declared value in the auction.

\subsection{Calibration of $\zeta$}
\label{sec:calibration}

We calibrate the externality cost $\zeta$ using two published estimates of the
dollar value of social media content externalities.

\paragraph{Impression-based estimate (Goldstein et al.).}
Taking a CPM (cost per thousand impressions) of \$1.53, the implied cost per unit
of aggregate rating is
\begin{equation}
    \zeta^{\text{imp}}
    = \frac{1.53/1000}{\;\overline{\beta_i/\eta_i}\;}
    \approx \$2.71 \text{ per rating unit per month,}
\end{equation}
where $\overline{\beta_i/\eta_i}$ is the sample mean of aggregate rating per impression.

\paragraph{Action-based estimate (Acquisti et al.).}
Taking a willingness-to-pay of \$0.54 per engagement action, the implied cost is
\begin{equation}
    \zeta^{\text{act}}
    = \frac{0.54}{\;\overline{\beta_i/\alpha_i}\;}
    \approx \$2.54 \text{ per rating unit per month,}
\end{equation}
where $\overline{\beta_i/\alpha_i}$ is the sample mean of aggregate rating per action.

Both methods yield $\zeta \approx \$2.5$, consistent with each other.  Because
there is inherent uncertainty in these estimates, we treat $\zeta$ as a free
parameter and report results across the range $\zeta \in [0.01,\, 100]$.


% ─────────────────────────────────────────────────────────────────────────────
\section{Auction Simulation}
\label{sec:simulation}
% ─────────────────────────────────────────────────────────────────────────────

\subsection{Setup}

For each experimental cell $(k, \zeta, D, n)$ — allocation slots, externality cost,
polynomial degree, and bidders per draw — we maintain two independent sets of auction
draws with different purposes.

\paragraph{Training draws ($J_{\text{train}} = 500$).}
These are used exclusively by the genetic algorithm as its fitness signal during
optimization.  The training seed is a function of $(\zeta, n)$ only — it excludes $k$
and $D$ — so all degrees and slot counts within the same $(\zeta, n)$ cell train on
identical post samples.  This ensures that any difference in optimized welfare across
degrees or $k$ values reflects the effect of those parameters and not sampling
variation in the training data.

\paragraph{Evaluation draws ($J_{\text{eval}} = 2000$).}
These are held out from the optimizer entirely and used only for final welfare
reporting.  The evaluation seed is offset from the training seed by a fixed constant,
guaranteeing independence between the two sets.  Like the training seed, the
evaluation seed excludes $k$ and $D$, so welfare numbers are directly comparable
across those dimensions for any fixed $(\zeta, n)$.  The larger draw count (2000
vs.\ 500) reduces Monte Carlo variance in the reported welfare estimates.

Each draw $j$ is formed by sampling $n$ posts without replacement from the dataset.
Post $i$ enters draw $j$ with pair $(v_i, e_i)$ computed
via~\eqref{eq:v_score} and~\eqref{eq:e_score}.

\subsection{Threshold Function}

The collateralized mechanisms use a polynomial threshold of degree $D$:
\begin{equation}
    \tau(e;\,\bsc) = \sum_{d=0}^{D} c_d\, e^d,
    \qquad \bsc = (c_0, c_1, \ldots, c_D) \in \mathbb{R}^{D+1}.
    \label{eq:tau}
\end{equation}
The coefficient vector $\bsc$ is the decision variable optimized by the genetic
algorithm (Section~\ref{sec:ga}).

\subsection{Three Auction Mechanisms}
\label{sec:mechanisms}

\subsubsection{VCG (Counterfactual)}

All $n$ advertisers participate.  The $k$ advertisers with the highest values
are selected:
\begin{equation}
    \mathcal{S}^{\text{vcg}}_j
    = \operatorname*{arg\,max}_{S \subseteq [n],\,|S|=k}\;\sum_{i \in S} v_i.
\end{equation}
Social welfare in draw $j$ is
\begin{equation}
    W^{\text{vcg}}_j = \sum_{i \in \mathcal{S}^{\text{vcg}}_j} (v_i + e_i).
\end{equation}

\subsubsection{Participant Audit (VCGA)}

Only advertisers whose declared value clears the threshold are admitted:
\begin{equation}
    \mathcal{M}_j(\bsc) = \bigl\{\, i \in [n] : v_i \geq \tau(e_i;\,\bsc) \bigr\}.
    \label{eq:admitted}
\end{equation}
Among admitted advertisers, the top $k$ by value are selected:
\begin{equation}
    \mathcal{S}^{\text{vcgPA}}_j(\bsc)
    = \operatorname*{arg\,max}_{S \subseteq \mathcal{M}_j,\,|S| \leq \min(k,\,|\mathcal{M}_j|)}
      \;\sum_{i \in S} v_i.
    \label{eq:vcga_winners}
\end{equation}
If fewer than $k$ advertisers are admitted, all admitted advertisers are selected.
Social welfare is
\begin{equation}
    W^{\text{vcgPA}}_j(\bsc) = \sum_{i \in \mathcal{S}^{\text{vcgPA}}_j} (v_i + e_i).
\end{equation}
The auction payment (generalised $(k{+}1)$th-price rule among admitted bidders) is
\begin{equation}
    p^{\text{PA}}_j(\bsc) =
    \begin{cases}
        v^{\mathcal{M}}_{(k+1)} & |\mathcal{M}_j| \geq k+1 \\
        0                        & \text{otherwise,}
    \end{cases}
    \label{eq:pa_payment}
\end{equation}
where $v^{\mathcal{M}}_{(k+1)}$ is the $(k{+}1)$th highest value among the admitted set.

\subsubsection{Recipient Audit (RA)}

No participation threshold is applied.  Each advertiser's individual social value
is $w_i = v_i + e_i$.  The $k$ advertisers with the largest \emph{positive} social
value are selected:
\begin{equation}
    \mathcal{S}^{\text{vcgRA}}_j
    = \operatorname*{arg\,max}_{S \subseteq \{\,i : w_i > 0\,\},\,|S| \leq k}
      \;\sum_{i \in S} w_i.
    \label{eq:ra_winners}
\end{equation}
If fewer than $k$ advertisers have positive social value, all such advertisers are
selected.  Social welfare is
\begin{equation}
    W^{\text{vcgRA}}_j = \sum_{i \in \mathcal{S}^{\text{vcgRA}}_j} w_i.
\end{equation}

\subsection{Welfare Aggregation}

All reported welfare numbers use the $J_{\text{eval}} = 2000$ evaluation draws.
Mean welfare for each mechanism is
\begin{equation}
    \bar{W}^{\text{vcg}} = \frac{1}{J_{\text{eval}}}\sum_{j=1}^{J_{\text{eval}}} W^{\text{vcg}}_j, \quad
    \bar{W}^{\text{vcgPA}}(\bsc^*) = \frac{1}{J_{\text{eval}}}\sum_{j=1}^{J_{\text{eval}}} W^{\text{vcgPA}}_j(\bsc^*), \quad
    \bar{W}^{\text{vcgRA}} = \frac{1}{J_{\text{eval}}}\sum_{j=1}^{J_{\text{eval}}} W^{\text{vcgRA}}_j,
\end{equation}
where $\bsc^*$ is the optimal threshold found by the genetic algorithm.
Welfare gains relative to the VCG counterfactual are
\begin{equation}
    \Delta W^{\text{PA}} = \bar{W}^{\text{vcgPA}}(\bsc^*) - \bar{W}^{\text{vcg}},
    \qquad
    \Delta W^{\text{RA}} = \bar{W}^{\text{vcgRA}} - \bar{W}^{\text{vcg}}.
    \label{eq:delta_w}
\end{equation}
Both quantities are reported in dollars per month per auction draw.


% ─────────────────────────────────────────────────────────────────────────────
\section{Genetic Algorithm Optimization}
\label{sec:ga}
% ─────────────────────────────────────────────────────────────────────────────

The goal is to find
\begin{equation}
    \bsc^* = \operatorname*{arg\,max}_{\bsc \in \mathcal{C}}\;
             \bar{W}^{\text{vcgPA}}(\bsc),
\end{equation}
where $\mathcal{C}$ is the gene space defined below.  Because the objective is
non-convex in $\bsc$ and the feasible space is continuous and high-dimensional,
we use a genetic algorithm (GA) operating directly in original $(v, e)$ space
without any normalization.

\subsection{Gene Space}
\label{sec:gene_space}

The gene space $\mathcal{C} = \mathcal{C}_0 \times \mathcal{C}_1 \times \cdots \times \mathcal{C}_D$
is derived from the empirical distribution of $(v_i, e_i)$ values so that the
search range scales correctly regardless of the externality cost $\zeta$ or the
polynomial degree $D$.

Let $\sigma_v$ denote the standard deviation of advertiser values across all
bidder draws, and let $\overline{|e|^d} = \mathbb{E}[|e_i|^d]$ denote the
mean absolute $d$th power of externalities.

The intercept gene space ensures the threshold can start above or below all bids:
\begin{equation}
    \mathcal{C}_0 = \bigl[v_{\min} - \tfrac{1}{2}\sigma_v,\;\;
                          v_{\max} + \tfrac{1}{2}\sigma_v\bigr].
    \label{eq:gene0}
\end{equation}

For slope coefficients ($d \geq 1$), the range is calibrated so that at a typical
externality value, the contribution $|c_d \cdot e^d|$ can shift $\tau$ by up to
$\lambda \sigma_v$ — covering the full spread of advertiser values:
\begin{equation}
    \mathcal{C}_d = \Bigl[{-\lambda\,\sigma_v}/{\overline{|e|^d}},\;\;
                           {\lambda\,\sigma_v}/{\overline{|e|^d}}\Bigr],
    \quad d = 1, \ldots, D,
    \label{eq:gened}
\end{equation}
where $\lambda = 3$ is the range multiplier (Table~\ref{tab:ga_hyperparams}).

\subsection{Vectorized Fitness Evaluation}

All $N_{\text{pop}}$ candidate solutions in the current population are evaluated
simultaneously using a single matrix multiply per auction draw.  Let
$\mathbf{P} \in \mathbb{R}^{N_{\text{pop}} \times (D+1)}$ be the population matrix
(row $g$ is candidate $\bsc^{(g)}$) and let
$\mathbf{E}^{(j)} \in \mathbb{R}^{n \times (D+1)}$ be the Vandermonde-style
power matrix for draw $j$ with entries $E^{(j)}_{ip} = e_i^{\,p}$.
Then
\begin{equation}
    \mathbf{T}^{(j)} = \mathbf{P}\,\bigl(\mathbf{E}^{(j)}\bigr)^{\!\top}
    \in \mathbb{R}^{N_{\text{pop}} \times n},
    \qquad T^{(j)}_{g,i} = \tau(e_i;\,\bsc^{(g)}),
    \label{eq:tau_mat}
\end{equation}
evaluating the threshold for every population member and every advertiser at once.
The admission mask is $A^{(j)}_{g,i} = \mathbf{1}[v_i \geq T^{(j)}_{g,i}]$.
Advertisers are sorted by $v$ descending; the cumulative sum of admitted indicators
identifies the first $k$ admitted positions, yielding the winner set.
Welfare is accumulated by summing $v_i + e_i$ over winners for all $g$ simultaneously.
The fitness of candidate $g$ is
\begin{equation}
    F_g = \frac{1}{J_{\text{train}}}\sum_{j=1}^{J_{\text{train}}} W^{\text{vcgPA}}_j(\bsc^{(g)}),
\end{equation}
accumulated across the $J_{\text{train}} = 500$ training draws without storing
per-draw matrices.  The evaluation draws are never seen by the optimizer.

\subsection{Selection, Crossover, and Mutation}

\paragraph{Selection.}
Tournament selection: for each parent slot, a random subset of the population is
drawn without replacement and the candidate with the highest fitness wins.

\paragraph{Crossover.}
Single-point crossover: a random cut point in $\{1,\ldots,D\}$ is chosen for each
pair of parents, and the offspring receives coefficients from parent~A before the
cut and from parent~B after.

\paragraph{Mutation.}
Gaussian perturbation with per-gene step sizes proportional to each gene's range:
\begin{equation}
    \sigma_d^{\text{mut}} = \sigma_{\text{frac}} \cdot |\mathcal{C}_d|,
    \label{eq:mut_sigma}
\end{equation}
where $|\mathcal{C}_d| = c_d^{\max} - c_d^{\min}$ is the gene range width and
$\sigma_{\text{frac}}$ is the mutation sigma fraction (Table~\ref{tab:ga_hyperparams}).
Each gene is independently perturbed with probability $p_{\text{mut}}$; to
prevent stagnation, at least one gene per offspring is always mutated.
All perturbed values are clipped to the gene bounds.

\subsection{Multi-Restart and Warm Start}
\label{sec:restarts}

To increase coverage of the coefficient space, each experimental cell is run for
$R$ independent restarts, each with a different random seed.  The solution with
the highest fitness across all restarts is retained.

\paragraph{Warm start.}
When sweeping polynomial degrees $D = 1, 2, 3$ in ascending order, the optimal
solution from degree $D$ is used to seed the first population member for
degree $D+1$:
\begin{equation}
    \bsc_{\text{seed}}^{(D+1)} = \bigl(c_0^*, c_1^*, \ldots, c_D^*,\; 0\bigr).
    \label{eq:warm_start}
\end{equation}
The warm start is applied only to restart~0; remaining restarts use a fresh
random population, preserving the benefit of the warm start while maintaining
diversity.

\paragraph{Fair cross-degree comparison.}
All degrees within the same cell $(k, \zeta, n)$ use an identical set of $J$
bidder draws (controlled by a shared sample seed that excludes the degree index).
This ensures that welfare differences across degrees reflect the effect of
polynomial flexibility rather than sampling variation.

\subsection{Hyperparameters}
\label{sec:hyperparams}

\begin{table}[H]
\caption{Genetic algorithm hyperparameters (defaults used unless noted otherwise).}
\label{tab:ga_hyperparams}
\begin{center}
\begin{tabularx}{\textwidth}{llX}
\toprule
Symbol & Default & Description \\
\midrule
$D$ & 1, 2, 3 & Polynomial degree of $\tau$ (swept) \\
$N_{\text{pop}}$ & 50 & Population size \\
$T$ & 500 & GA generations per restart \\
$R$ & 1 & Independent restarts per cell \\
$N_{\text{mate}}$ & 15 & Parents selected for crossover each generation \\
$p_{\text{mut}}$ & 0.25 & Per-gene mutation probability \\
$\sigma_{\text{frac}}$ & 0.15 & Gaussian mutation step as fraction of gene range \\
$\lambda$ & 3.0 & Gene range multiplier (Eq.~\ref{eq:gened}) \\
$J_{\text{train}}$ & 500  & Training auction draws per cell (GA fitness) \\
$J_{\text{eval}}$  & 2000 & Evaluation draws per cell (welfare reporting only) \\
$n$ & 20 & Bidders sampled per auction draw \\
\bottomrule
\end{tabularx}
\end{center}
\end{table}


% ─────────────────────────────────────────────────────────────────────────────
\section{Expected Penalty Curves}
\label{sec:penalties}
% ─────────────────────────────────────────────────────────────────────────────

Given the optimized threshold $\tau(e;\,\bsc^*)$, we compute the expected penalty
an advertiser with externality value $e$ would face under each audit mechanism.

\subsection{Participant Audit: $R^{\text{PA}}(e)$}
\label{sec:r_pa}

Under the participant audit, an advertiser with externality $e$ is excluded from
any draw in which the market-clearing price $p^{\text{PA}}_j$ falls below the
threshold $\tau(e;\,\bsc^*)$.  To clear the threshold in such a draw, the advertiser
would need to post collateral equal to the shortfall
$\tau(e;\,\bsc^*) - p^{\text{PA}}_j > 0$.

Let $\{p_j\}_{j=1}^{J_{\text{eval}}}$ be the sequence of auction payments from
\eqref{eq:pa_payment}, computed on the evaluation draws.  Define
\begin{align}
    A(e) &= \frac{1}{J_{\text{eval}}}\bigl|\bigl\{ j : p_j < \tau(e;\,\bsc^*) \bigr\}\bigr|,
    \label{eq:A_e} \\
    P(e) &= \mathbb{E}\!\left[p_j \;\big|\; p_j < \tau(e;\,\bsc^*)\right].
    \label{eq:P_e}
\end{align}
Here $A(e)$ is the fraction of draws in which the market price is below the
threshold, and $P(e)$ is the mean price conditional on being below the threshold.
The expected participant penalty at externality $e$ is
\begin{equation}
    R^{\text{PA}}(e) = A(e)\,\bigl[\tau(e;\,\bsc^*) - P(e)\bigr].
    \label{eq:r_pa}
\end{equation}
For large negative externalities, the threshold $\tau(e;\bsc^*)$ is high relative
to the market price, so both $A(e)$ and the expected gap are large, resulting in a
steep penalty.  For posts with small or positive externalities, $\tau(e;\bsc^*)$ is
low and most market prices clear it, so $A(e) \approx 0$ and $R^{\text{PA}}(e)
\approx 0$.

\subsection{Recipient Audit: $R^{\text{RA}}(e)$}
\label{sec:r_ra}

Under the recipient audit, selection is based on social value $w_i = v_i + e_i$.
An advertiser with externality $e$ internalises the full cost of its externality:
\begin{equation}
    R^{\text{RA}}(e) = -e.
    \label{eq:r_ra}
\end{equation}
This is linear in $e$: advertisers generating negative externalities ($e < 0$)
pay a penalty equal to $|e|$, while those generating positive externalities
($e > 0$) effectively receive a subsidy.

\end{document}
